\documentclass[tikz,border=8pt]{standalone}
\usepackage{ctex}
\usepackage{amsmath}
\usepackage{pgfplots}
\pgfplotsset{compat=1.18}
\usetikzlibrary{calc,positioning,arrows.meta}

%% 图 opt-p2-06-2：凸 vs 非凸优化保证
%% 左：凸函数（GD 全局收敛，唯一极小）
%% 右：非凸函数（GD 收敛到局部极小，全局极小未必到达）

\begin{document}
\begin{tikzpicture}

%% ===== 左图：凸函数 =====
\begin{scope}[xshift=0cm]

  % 坐标轴
  \draw[->, gray!55, thin] (-2.5, -0.3) -- (2.5, -0.3)
    node[right, font=\scriptsize] {$x$};
  \draw[->, gray!55, thin] (-2.5, -0.3) -- (-2.5, 3.5)
    node[above, font=\scriptsize] {$f(x)$};

  % 凸函数：f(x) = x^2 + 0.3（抛物线，唯一极小）
  \draw[blue!80!black, very thick, domain=-1.9:1.9, samples=60]
    plot (\x+0, {\x*\x + 0.3 - 0.3}) node[right, font=\scriptsize] {};

  % 唯一全局极小点
  \fill[red!80!black] (0, 0) circle (3pt)
    node[below, font=\small, yshift=-2pt] {$\mathbf{x}^*$（全局极小）};

  % 三条从不同初始点出发的梯度下降轨迹
  % 全部收敛到 x* = 0
  % 轨迹1：从 x=-1.8 出发
  \draw[->, green!60!black, thick, >=Stealth]
    (-1.8, 3.24) -- (-1.0, 1.0);
  \draw[->, green!60!black, thick, >=Stealth]
    (-1.0, 1.0) -- (-0.35, 0.12);
  \draw[->, green!60!black, thick, dashed, >=Stealth]
    (-0.35, 0.12) -- (-0.05, 0.003);
  \fill[green!60!black] (-1.8, 3.24) circle (2.5pt)
    node[above left, font=\scriptsize] {$x_0^{(1)}$};

  % 轨迹2：从 x=1.6 出发
  \draw[->, orange!70!black, thick, >=Stealth]
    (1.6, 2.56) -- (0.8, 0.64);
  \draw[->, orange!70!black, thick, >=Stealth]
    (0.8, 0.64) -- (0.25, 0.063);
  \draw[->, orange!70!black, thick, dashed, >=Stealth]
    (0.25, 0.063) -- (0.04, 0.002);
  \fill[orange!70!black] (1.6, 2.56) circle (2.5pt)
    node[above right, font=\scriptsize] {$x_0^{(2)}$};

  % "全局收敛" 标注
  \node[font=\small, green!50!black, align=center]
    at (0, 3.6) {所有初始点\\均收敛到 $\mathbf{x}^*$};
  \draw[->, green!50!black, thin] (0, 3.3) -- (0, 0.15);

  % 下方标注
  \node[font=\small\bfseries, blue!80!black] at (0, -0.9) {凸函数};
  \node[font=\scriptsize, align=center] at (0, -1.35)
    {$\nabla^2 f \succeq 0$\\局部极小 $=$ 全局极小\\一阶条件充要};

\end{scope}

%% ===== 右图：非凸函数 =====
\begin{scope}[xshift=7.0cm]

  % 坐标轴
  \draw[->, gray!55, thin] (-2.5, -0.3) -- (2.5, -0.3)
    node[right, font=\scriptsize] {$x$};
  \draw[->, gray!55, thin] (-2.5, -0.3) -- (-2.5, 3.8)
    node[above, font=\scriptsize] {$f(x)$};

  % 非凸函数：f(x) = x^4 - 2.5*x^2 + 2（双极小 + 一极大）
  % 极小点约 x = ±1.12，极大点 x = 0
  % f(x) = 0.5*(x^4 - 3*x^2) + 2.5
  \draw[purple!70!black, very thick, domain=-1.95:1.95, samples=80]
    plot (\x, {0.5*(\x*\x*\x*\x - 3*\x*\x) + 2.5});

  % 两个局部（全局）极小点 x ≈ ±1.22, f ≈ 1.1
  \fill[red!80!black] (-1.22, 1.1) circle (3pt)
    node[below, font=\scriptsize, yshift=-2pt] {局部极小};
  \fill[red!20!black] (1.22, 1.1) circle (3pt)
    node[below, font=\scriptsize, yshift=-2pt] {全局极小};
  \node[font=\scriptsize, red!20!black] at (1.22, 0.5) {$\mathbf{x}^{**}$};
  \node[font=\scriptsize, red!80!black] at (-1.22, 0.5) {$\tilde{\mathbf{x}}$};

  % 鞍点（极大点）x=0
  \fill[orange!80!black] (0, 2.5) circle (2.5pt)
    node[above right, font=\scriptsize] {鞍点};

  % 轨迹1：从左侧出发，收敛到局部极小（非全局）
  \draw[->, red!80!black, thick, >=Stealth]
    (-1.8, 3.1) -- (-1.5, 1.6);
  \draw[->, red!80!black, thick, >=Stealth]
    (-1.5, 1.6) -- (-1.28, 1.12);
  \draw[->, red!80!black, thick, dashed, >=Stealth]
    (-1.28, 1.12) -- (-1.23, 1.1);
  \fill[red!80!black] (-1.8, 3.1) circle (2.5pt)
    node[above left, font=\scriptsize] {$x_0^{(1)}$};

  % 轨迹2：从右侧出发，收敛到全局极小
  \draw[->, green!60!black, thick, >=Stealth]
    (1.8, 3.1) -- (1.5, 1.6);
  \draw[->, green!60!black, thick, >=Stealth]
    (1.5, 1.6) -- (1.28, 1.12);
  \draw[->, green!60!black, thick, dashed, >=Stealth]
    (1.28, 1.12) -- (1.23, 1.1);
  \fill[green!60!black] (1.8, 3.1) circle (2.5pt)
    node[above right, font=\scriptsize] {$x_0^{(2)}$};

  % 警示标注：依赖初始点
  \node[font=\scriptsize, red!70!black, fill=yellow!20,
        draw=red!40, thin, rounded corners=2pt, inner sep=3pt, align=center]
    at (-1.22, 3.0) {收敛到\\局部极小！};
  \draw[->, red!60!black, thin] (-1.22, 2.7) -- (-1.22, 1.6);

  % 下方标注
  \node[font=\small\bfseries, purple!70!black] at (0, -0.9) {非凸函数};
  \node[font=\scriptsize, align=center] at (0, -1.35)
    {$\nabla^2 f$ 不定（含鞍点）\\GD 收敛取决于初始点\\全局极小无保证};

\end{scope}

%% ===== 整体标题 =====
\node[font=\large\bfseries] at (6.0, 4.5)
  {凸 vs 非凸优化：全局收敛保证的差异};

%% ===== 对比总结框 =====
\node[draw=gray!50, thin, fill=gray!5, rounded corners=3pt,
      font=\scriptsize, inner sep=5pt, align=left]
  at (6.0, -2.2)
  {凸：$\nabla f=0 \;\Leftrightarrow\;$ 全局极小（充要），GD 总收敛 \quad
   非凸：$\nabla f=0$ 只是必要，GD 可能停在局部极小或鞍点};

\end{tikzpicture}
\end{document}
